\section{Arbre de diffusion}

On suppose qu'on dispose d'un arbre couvrant $T$ enraciné en $r_0$ donné par $\begin{cases*}u \mapsto \mr{parent}(u) \\ u \mapsto \mr{Enfants}(u) \end{cases*}$
On considre l'algorithme de diffusion suivant~:

\begin{lstLNat}
    CAST(T, r0):
      Si $u \neq r_0$: 
        M := RECEIVE()
      Pour tout $v$ dans $\mr{Enfants}(u)$:
        SEND(M,v)
\end{lstLNat}

\begin{proposition}{}{propriétés de cet algorithme}
  \begin{itemize}
    \item $\mr{TEMPS}\Big(CAST(T,r_0),G\Big) = \mr{ecc}_T(r_0) = \mr{hauteur}(T)$
    \item $\mr{MESSAGE}\Big(CAST(T,r_0),G\Big) = \abs{E(T)} = n-1$
  \end{itemize}
\end{proposition}

Cet algorithme a le mérite de limiter le nombre de messages à $n-1$ mais il nécessite d'avoir $T$. 

Considérons à présent l'algorithme asynchrone suivant~:
\begin{lstLNat}
  FLOOD(T, r0):
    si $u = r_0$:
      SENDALL(M) // envoi à tous les voisins
    sinon:
      M := receive()
      v := émetteru du message M
      SEND(M, w) pour tout $w \in \mr{Voisins}(u) \setminus \{v$
\end{lstLNat}
\begin{proposition}{}{propriétés de cet algorithme}
  On a, en notant $m$ le nombre d'arêtes de $G$, et $n$ son nombre de sommets~:
  \begin{itemize}
    \item $\mr{MESSAGE}\Big(FLOOD(r_0),G\Big) = m = \bigO(n^2)$
    \item $\mr{TEMPS}\Big(FLOOD(r_0),G\Big) = \mr{ecc}_G(r_0)$
  \end{itemize}
\end{proposition}
\begin{demonstration}
  \begin{enumerate}
    \item\begin{align*}
        \mr{MESSAGE}\Big(FLOOD(r_0),G\Big) &= \deg_G(r_0) \sum_{\substack{u \in V\\ u \neq r_0}}\Big(\mr{deg}_G(u) - 1\Big) \\
        &= \sum_{u \in V} \deg(u) - (n-1) \\
        &= 2m - n +1
      \end{align*}
    \item Montre qu'après $t$ unités de temps, tous les sommets à distances $\leq t$ ont reçu $M$. Il y a donc au plus $\mr{ecc}_G(r_0)$ unités de temps avant que tous les sommets aient reçu $M$. 
  \end{enumerate}
\end{demonstration}


Il existe également une variante synchrone "amnésique" de FLOOD qui consiste à appliquer FLOOD perpétuellement.
\begin{lstLNat}
  AMNESIAC-FLOOD(r0):
  si $u=r_0$:
    SENDALL(M)
  S = $\varnothing$
  tant que vrai:
    démarrer une nouvelle ronde
    si $S = \varnothing$:
      SEND(M, v) pour tout $v \in \mr{Voisins}(u)\setminus S$
      // c'est grâce à ça qu'on s'arrête
      $\mc{M} := RECEIVE_ALL()$ // on reçoit tous les messages des voisins
      $S$ := l'ensemble des émetteurs de $\mc{M}$
      si $S \neq \varnothing$:
        choisir $M \in \mc{M}$
\end{lstLNat}

\subsection{Message \textit{vs.} connaissance}

\begin{tabular}{c|c|c|c}
    & TEMPS & MESSAGE & "CONNAISSANCE" \\
    \hline
    $\mr{CAST}_T$ & $\mr{ecc}_T(r_0)$ & $n-1 = \bigO(n)$ & $T$ \\
    \hline
    $\mr{FLOOD}$ & $\mr{ecc}_G(r_0)$ & $2m-n+1 = \bigO(n^2)$ & $\varnothing$ \\
    \hline
\end{tabular}

\begin{definition}{}{connaissance à distance $\rho$}
  On dit qu'un sommet $u$ a connaissance à distance $\rho$ si $u$ connaît l'ensemble des sommets à distance $\leq \rho$ de lui, ainsi que les arêtes entre ces sommets.
\end{definition}

\begin{proposition}{}{petits messages}
  $\forall \rho \in \N, \forall G$ connexe à $n$ sommets et $r_0$, il existe un algorithme de diffusion $\mr{FLOODK}_\rho$ avec une connaissance $\rho$ tel que~:
  $$\mr{MESSAGE}\Big(\mr{FLOODK}_\rho(r_0), G\Big) = 2\min(m, n^{1 + \frac{1}{\rho}})$$
\end{proposition}


Ainsi~:
\begin{itemize}
  \item $\mr{FLOODK}_0$ est en $\bigO(n^2)$ 
  \item $\mr{FLOODK}_1$ est l'algorithme $\mr{FLOOD}$, également en $\bigO(n^2)$ qui ne nécessite aucune connaissance.
  \item $\mr{FLOODK}_2$ est en $\bigO(n^{3/2})$
  \item $\mr{FLOODK}_3$ est en $\bigO(n^{4/3})$
  \item $\mr{FLOODK}_\infty$ est l'algorithme $\mr{CAST}_T$, en $\bigO(n)$
\end{itemize}


\begin{demonstration}
  On prend $\rho \geq 1$. Deux idées~:
\begin{itemize}
  \item Localement, chaque sommet décide de désactiver certaines arêtes, et $\tilde{G}$ le graphe résultant.
  \item On applique FLOOD dans $\tilde{G}$, $\mr{MSG} \leq 2\tilde{m} - {n} + 1$
\end{itemize}
\end{demonstration}

\newcommand{\id}{\mr{ID}}

On associe un identifiant $\id(u,v)$ à chaque arête $(u,v)$. On impose la contrainte suivante~:
\begin{proposition}{}{}
  $u$ désactive $u-v$ s'il existe dans $G$ un cycle $\mc{C}$ de longueur $\leq 2\rho$ contenant $u-v$, d'identifiant minimal dans ce cycle.
\end{proposition}
On parle ici de cycles au sens où l'on ne passe pas plusieurs fois par la même arête.
\begin{remarque}{}{}
  Chaque sommet peut appliquer la règle localement, mais pas sur le même cycle. On ne peut pas être facilement sûr que le graphe d'arrivée soit connexe, sachant que toutes les arêtes sont enlevées non pas successivement, mais simultanément.
\end{remarque}

\begin{theoreme}{}{}
  Si $G$ est connexe, alors $\tilde{G}$ est connexe.
\end{theoreme}

\begin{demonstration}
  On montre le résultat par récurrence sur l'indice de l'arête supprimée (on les ordonne arbitrairement). On note $G_i$ le graphe après suppression des $i$ premières arêtes. On suppose que $G_i$ est connexe et on montre que $G_{i+1}$ est connexe. 

  On note $e = (u,v)$ l'arête supprimée à l'étape $i+1$. Par hypothèse, il existe un cycle $\mc{C}$ de longueur $\leq 2\rho$ contenant $e$. Comme $\mc{C}$ est un cycle, il existe un chemin de $u$ à $v$ dans $\mc{C} \setminus \{e\}$. Ce chemin est contenu dans $G_i$, donc $u$ et $v$ sont encore connectés dans $G_{i+1}$. Ainsi, la suppression de l'arête ne déconnecte pas le graphe.
\end{demonstration}

On peut se questionner sur le majorant $2 \rho$. Il découle du résultat suivant~:

\begin{theoreme}{}{}
  Si $G$ n'a pas de cycle de longueur $\leq 2\rho$, alors il a au plus $n^{1+\frac{1}{\rho}} + n$
\end{theoreme}


\begin{demonstration}
  On note $d = \frac{m}{n}$ le degré moyen. On supprime tous les sommets de degré $<d$ de $G$. Ainsi, on obtient un sous-graphe $H$ de $G$, qui contient au moins un arbre $\mc{A}$ complet d'arité $d$, et de profondeur $\rho$. On compte les feuilles de cet arbre $\mc{A}$~: il y en a $(d-1)^\rho$. Or, ces feuilles sont des sommets de $H$, donc de $G$. On a alors~:
  $$ n \geq (d-1)^\rho $$
Donc~: 
$$ n \geq \left(\frac{m}{n} - 1\right)^\rho $$
D'où le résultat.
\end{demonstration}
